\documentclass{../../nndlcourse}

\begin{document}

\title{实验3 - 卷积神经网络}
\author{数据科学与大数据技术专业}
\date{第6周}
\maketitle


% ============================================================
\section{实验目标}
% ============================================================
\begin{itemize}
    \item 理解卷积神经网络(CNN)中局部连接、权值共享、感受野、填充、步幅和池化的作用.
    \item 使用 NumPy 从零实现卷积层和池化层的前向传播, 并理解输出尺寸的计算方式.
    \item 选做实现卷积层和池化层的反向传播, 理解梯度如何传回输入、滤波器和偏置.
    \item 使用 PyTorch 构建卷积神经网络, 完成 Happy House 二分类和 SIGNS 手语数字多分类任务.
    \item 掌握 \texttt{nn.Sequential} 与自定义 \texttt{nn.Module} 两种模型组织方式的差异.
    \item 熟悉图像张量在 NumPy 与 PyTorch 中的形状转换, 尤其是 \texttt{NHWC} 与 \texttt{NCHW} 格式.
\end{itemize}


% ============================================================
\section{实验环境}
% ============================================================
\begin{itemize}
    \item 操作系统: Windows / Linux / macOS
    \item Python 3.7 及以上
    \item 主要依赖库: \texttt{numpy}、\texttt{matplotlib}、\texttt{h5py}、\texttt{torch}
    \item 推荐开发环境: Jupyter Notebook 或 VS Code
    \item 实验文件:
    \begin{itemize}
        \item \texttt{W6A1/Convolution\_model\_Step\_by\_Step\_v1.ipynb}
        \item \texttt{W6A2/Convolution\_model\_Application\_PyTorch.ipynb}
    \end{itemize}
\end{itemize}

如本地环境缺少依赖, 可在终端中执行:
\begin{verbatim}
pip install numpy matplotlib h5py torch
\end{verbatim}


% ============================================================
\section{背景知识}
% ============================================================

\subsection{卷积神经网络的基本思想}

全连接网络在处理图像时会把图像展平为长向量, 这会破坏像素的空间邻接结构, 并导致参数量迅速增大. 卷积神经网络通过局部连接和权值共享来直接处理图像张量, 使模型能够学习边缘、纹理、局部形状等空间模式.

对输入图像体积 $A_{\mathrm{prev}}\in\mathbb{R}^{n_H\times n_W\times n_C}$, 一个卷积滤波器 $W_c\in\mathbb{R}^{f\times f\times n_C}$ 在某个空间位置上计算:
\begin{align}
    z = \sum_{i=1}^{f}\sum_{j=1}^{f}\sum_{k=1}^{n_C}
    A_{\mathrm{slice}, i,j,k} W_{c,i,j,k} + b_c
\end{align}
同一个滤波器会滑过整张图像, 产生一个二维特征图. 多个滤波器堆叠后形成多通道输出.

\subsection{填充、步幅与输出尺寸}

\textbf{填充(Padding)}是在图像边界补零. 它可以保留边缘信息, 也可以控制卷积后的空间尺寸. \textbf{步幅(Stride)}表示滑动窗口每次移动的像素数.

若输入高度和宽度分别为 $n_{H_{\mathrm{prev}}}$、$n_{W_{\mathrm{prev}}}$, 滤波器大小为 $f$, 填充为 $p$, 步幅为 $s$, 则卷积输出尺寸为:
\begin{align}
    n_H &= \left\lfloor \frac{n_{H_{\mathrm{prev}}}-f+2p}{s}\right\rfloor + 1 \\
    n_W &= \left\lfloor \frac{n_{W_{\mathrm{prev}}}-f+2p}{s}\right\rfloor + 1 \\
    n_C &= \text{滤波器个数}
\end{align}
本实验第一部分的数据格式采用 NumPy 常见的 \texttt{NHWC}: $(m,n_H,n_W,n_C)$.

\subsection{池化层}

池化层没有可学习参数, 主要用于降低空间尺寸和增强局部平移不变性. 最大池化取窗口内最大值:
\begin{align}
    A_{h,w,c} = \max A_{\mathrm{prev}}[h:h+f,\,w:w+f,\,c]
\end{align}
平均池化取窗口均值:
\begin{align}
    A_{h,w,c} = \frac{1}{f^2}\sum A_{\mathrm{prev}}[h:h+f,\,w:w+f,\,c]
\end{align}
若池化窗口大小为 $f$, 步幅为 $s$, 且不使用填充, 输出尺寸为:
\begin{align}
    n_H = \left\lfloor \frac{n_{H_{\mathrm{prev}}}-f}{s}\right\rfloor + 1,\quad
    n_W = \left\lfloor \frac{n_{W_{\mathrm{prev}}}-f}{s}\right\rfloor + 1,\quad
    n_C = n_{C_{\mathrm{prev}}}
\end{align}

\subsection{卷积与池化的反向传播}

卷积层反向传播需要把输出梯度 $dZ$ 分别累加到输入、权重和偏置上. 对某一输出位置和滤波器 $c$:
\begin{align}
    dA_{\mathrm{prev,slice}} &\mathrel{+}= W_c\, dZ_{h,w,c} \\
    dW_c &\mathrel{+}= A_{\mathrm{prev,slice}}\, dZ_{h,w,c} \\
    db_c &\mathrel{+}= dZ_{h,w,c}
\end{align}
最大池化反向传播只把梯度传给窗口中最大值所在的位置; 平均池化反向传播则把梯度均匀分配给窗口内所有元素.

\subsection{PyTorch 中的图像张量格式}

PyTorch 的卷积层 \texttt{nn.Conv2d} 接收 \texttt{NCHW} 格式, 即 $(N,C,H,W)$. 而本实验数据集通常由 H5 文件读入为 \texttt{NHWC}, 即 $(N,H,W,C)$. 因此输入 PyTorch 前需要转置:
\begin{verbatim}
X_t = torch.tensor(X_np.transpose(0, 3, 1, 2), dtype=torch.float32)
\end{verbatim}
\begin{itemize}
    \item 二分类任务可使用 Sigmoid 输出并配合 \texttt{BCELoss}.
    \item 多分类任务通常直接输出 logits, 并使用交叉熵损失; 该损失内部已经包含 log softmax.
\end{itemize}


% ============================================================
\section{实验步骤}
% ============================================================

\subsection*{实验3.1: NumPy 实现卷积网络基础模块(W6A1)}

\textbf{实验目标: }不依赖深度学习框架, 用 NumPy 实现 CNN 的基础运算, 明确每个循环变量和张量形状的含义.

\subsubsection*{3.1.1 零填充}

\textbf{练习 1: \texttt{zero\_pad}}

对一批图像 $X\in\mathbb{R}^{m\times n_H\times n_W\times n_C}$ 在高度和宽度维度补零:
\begin{align}
    X_{\mathrm{pad}}\in\mathbb{R}^{m\times(n_H+2p)\times(n_W+2p)\times n_C}
\end{align}
实现要点:
\begin{itemize}
    \item 使用 \texttt{np.pad}.
    \item 只填充第 2、3 个维度, 不填充样本维和通道维.
    \item 推荐写法: \texttt{((0,0),(pad,pad),(pad,pad),(0,0))}.
\end{itemize}

\subsubsection*{3.1.2 单步卷积}

\textbf{练习 2: \texttt{conv\_single\_step}}

实现单个窗口和单个滤波器的卷积:
\begin{align}
    Z = \sum (a_{\mathrm{slice}}\odot W) + b
\end{align}
其中 $a_{\mathrm{slice}}$ 与 $W$ 形状均为 $(f,f,n_C)$. 实现时先逐元素乘法, 再 \texttt{np.sum}, 最后加上偏置并转为标量.

\subsubsection*{3.1.3 卷积前向传播}

\textbf{练习 3: \texttt{conv\_forward}}

将单步卷积推广到整个 batch、所有空间位置和所有滤波器. 核心循环为:
\begin{itemize}
    \item 遍历样本 $i$.
    \item 遍历输出高度 $h$ 与宽度 $w$.
    \item 根据 \texttt{stride} 计算 \texttt{vert\_start}、\texttt{vert\_end}、\texttt{horiz\_start}、\texttt{horiz\_end}.
    \item 遍历滤波器通道 $c$, 调用 \texttt{conv\_single\_step}.
\end{itemize}
函数返回卷积输出 \texttt{Z} 和反向传播需要的 \texttt{cache=(A\_prev,W,b,hparameters)}.

\subsubsection*{3.1.4 池化前向传播}

\textbf{练习 4: \texttt{pool\_forward}}

实现池化层的前向传播, 支持:
\begin{itemize}
    \item \texttt{mode="max"}: 最大池化.
    \item \texttt{mode="average"}: 平均池化.
\end{itemize}
池化层对每个通道独立操作, 不改变通道数. 实现时要注意池化窗口只覆盖高度和宽度, 不跨通道.

\subsubsection*{3.1.5 反向传播(选做)}

选做练习包括:
\begin{itemize}
    \item \texttt{conv\_backward}: 计算 \texttt{dA\_prev}、\texttt{dW}、\texttt{db}.
    \item \texttt{create\_mask\_from\_window}: 为最大池化反向传播生成最大值掩码.
    \item \texttt{distribute\_value}: 为平均池化反向传播均匀分配梯度.
    \item \texttt{pool\_backward}: 整合最大池化与平均池化的反向传播.
\end{itemize}
完成选做时要特别处理 \texttt{pad=0} 的情况, 避免用空切片覆盖 \texttt{dA\_prev}.

\subsection*{实验3.2: PyTorch 卷积网络应用(W6A2)}

\textbf{实验目标: }使用 PyTorch 搭建、训练和评估卷积神经网络, 比较 \texttt{nn.Sequential} 与自定义 \texttt{nn.Module}.

\subsubsection*{3.2.1 Happy House 二分类}

Happy House 使用 $64\times64$ 的 RGB 图像, 标签表示图片中的人是否微笑.
模型使用顺序容器构建:
\begin{align}
    \mathrm{ZeroPad2d}\to \mathrm{Conv2d}\to \mathrm{BatchNorm2d}\to
    \mathrm{ReLU}\to \mathrm{MaxPool2d}\to \mathrm{Flatten}\to
    \mathrm{Linear}\to \mathrm{Sigmoid}
\end{align}
关键形状推导:
\begin{itemize}
    \item 输入: $(N,3,64,64)$.
    \item \texttt{ZeroPad2d(3)} 后: $(N,3,70,70)$.
    \item \texttt{Conv2d(3,32,kernel\_size=7,stride=1)} 后: $(N,32,64,64)$.
    \item \texttt{MaxPool2d(2,2)} 后: $(N,32,32,32)$.
    \item 展平维度: $32\times32\times32=32768$.
\end{itemize}
训练时使用 \texttt{nn.BCELoss()} 和 Adam 优化器. 预测时以 0.5 为阈值计算二分类准确率.

\subsubsection*{3.2.2 SIGNS 手语数字多分类}

第二个任务是识别 0 到 5 的手语数字. 使用自定义 PyTorch 模块, 主体结构为两个``卷积 + ReLU + 池化''模块, 接展平层和线性分类层.
推荐结构:
\begin{itemize}
    \item \texttt{Conv2d(3,8,kernel\_size=4,stride=1,padding='same')}.
    \item \texttt{MaxPool2d(kernel\_size=8,stride=8)}, 空间尺寸 $64\to8$.
    \item \texttt{Conv2d(8,16,kernel\_size=2,stride=1,padding='same')}.
    \item \texttt{MaxPool2d(kernel\_size=4,stride=4)}, 空间尺寸 $8\to2$.
    \item \texttt{Flatten} 后维度为 $16\times2\times2=64$.
    \item \texttt{Linear(64,6)} 输出 6 类 logits.
\end{itemize}
多分类训练使用 \texttt{nn.CrossEntropyLoss()}. 标签应为整数类别索引, 不要转为 one-hot; 模型最后也不要加 softmax.

\subsubsection*{3.2.3 训练循环与结果分析}

PyTorch 标准训练循环包括:
\begin{enumerate}
    \item \texttt{optimizer.zero\_grad()} 清空梯度.
    \item 前向传播得到输出或 logits.
    \item 计算损失.
    \item \texttt{loss.backward()} 自动反向传播.
    \item \texttt{optimizer.step()} 更新参数.
\end{enumerate}
请记录训练集和测试集的损失、准确率, 并绘制训练历史曲线. 若训练准确率显著高于测试准确率, 应分析是否过拟合.


% ============================================================
\section{核心函数与模块总结}
% ============================================================

\begin{center}
\begin{tabular}{|l|l|l|}
\hline
\textbf{函数 / 模型} & \textbf{所在实验} & \textbf{作用} \\
\hline
\texttt{zero\_pad} & W6A1 & 对图像 batch 做零填充 \\
\texttt{conv\_single\_step} & W6A1 & 单窗口单滤波器卷积 \\
\texttt{conv\_forward} & W6A1 & 卷积层前向传播 \\
\texttt{pool\_forward} & W6A1 & 最大 / 平均池化前向传播 \\
\texttt{conv\_backward} & W6A1 选做 & 卷积层反向传播 \\
\texttt{pool\_backward} & W6A1 选做 & 池化层反向传播 \\
\texttt{happyModel} & W6A2 & \texttt{nn.Sequential} 二分类 CNN \\
\texttt{ConvolutionalModel} & W6A2 & 自定义 \texttt{nn.Module} 多分类 CNN \\
\hline
\end{tabular}
\end{center}


% ============================================================
\section{实验作业}
% ============================================================
\begin{enumerate}
    \item 完成 \texttt{W6A1} 中 \texttt{zero\_pad}、\texttt{conv\_single\_step}、\texttt{conv\_forward}、\texttt{pool\_forward} 的代码, 并通过测试单元格.
    \item \textbf{选做}: 完成 \texttt{conv\_backward}、\texttt{create\_mask\_from\_window}、\texttt{distribute\_value}、\texttt{pool\_backward}.
    \item 完成 \texttt{W6A2} 中 HappyModel 与 SIGNS 多分类模型的搭建、训练和评估.
    \item 在实验记录中说明 \texttt{NHWC} 到 \texttt{NCHW} 的转换原因, 并给出两个任务的损失曲线和准确率.
    \item 比较 \texttt{nn.Sequential} 与自定义 \texttt{nn.Module} 的适用场景.
    \item \textbf{选做}: 寻找一个小型 CNN 开源项目或公开教程, 在本地复现并用文字说明模型结构、数据集和运行结果.
\end{enumerate}


% ============================================================
\section{提交要求}
% ============================================================

请在提交前检查材料是否齐全. 本课程作业上传只接受文本文件, 不接受 PDF、Word 文档、图片、截图、压缩包等二进制文件. 采用这一规则, 一方面是为了引入 AI 辅助评阅, 另一方面也是希望大家养成用清晰文本表达实验过程和结论的习惯.

如果实验报告确实需要配图, 请优先思考是否可以用文字、公式或表格表达清楚; 若必须插图, 请将图片放入 Notebook 中, 形成一份完整的 \texttt{.ipynb} 实验报告. 若使用 \LaTeX{} 或 Word 撰写报告, 请使用 \texttt{pandoc} 等工具转换为 \texttt{.md} 或 \texttt{.txt} 等文本格式后再提交, 不要提交 \texttt{.pdf} 或 \texttt{.docx}.

\begin{enumerate}
    \item \textbf{Jupyter Notebook 文件(.ipynb)}
    \begin{enumerate}
        \item 必交内容包括 \texttt{W6A1} 和 \texttt{W6A2} 对应的 Notebook 文件.
        \item Notebook 中应包含完整代码、运行结果、损失曲线、准确率输出和必要的分析文字.
        \item 只需提交 Notebook 源文件, 不需要额外导出为 PDF、HTML、Python 脚本、Word 文档、截图或图片文件.
    \end{enumerate}
    \item \textbf{实验过程与 AI 互动记录文本文件(.md / .txt)}
    \begin{enumerate}
        \item 必须额外提交一份文本文件, 作为实验过程与 AI 互动记录, 推荐使用 \texttt{.md} 或 \texttt{.txt}.
        \item 记录内容至少包括: 查阅了哪些资料、向 AI 提出了哪些问题、如何继续思考、调试、修改代码或完善实验分析.
        \item 不要把 AI 的回复原文放入记录中, 也不要复制粘贴大段代码; 记录应只包含你自己写的内容, 包括你的问题、判断、尝试、修改思路和必要的简短总结.
        \item 记录应尽量按时间顺序整理, 不要只保留最终答案.
    \end{enumerate}
\end{enumerate}


% ============================================================
\section{关于作业的重要说明}
% ============================================================

本实验的核心目标是理解卷积和池化的底层计算, 以及 PyTorch 如何把这些计算封装为可训练模块. 代码填空部分应由学生自己完成, 并通过调试和单元测试逐步确认正确性.

可以使用 AI 工具辅助理解概念、解释报错、查询 API 用法或检查张量形状, 但不得直接让 AI 代写全部必做代码后不加理解地提交. 若使用 AI 工具, 必须在文本记录中如实写明自己的提问、自己的判断和后续修改过程.

AI 互动记录应只保留学生自己写的内容. 不要粘贴 AI 的完整回复, 不要粘贴复制来的大段代码, 也不要把聊天记录原样导出后提交.

选做内容可以更自由地使用 AI 工具, 不需记录使用过程。

\end{document}
